\section{西部边疆：王室为什么越来越难守住关中}

西周晚期的威胁不只来自内部。玁狁、戎狄等名称在不同材料中覆盖的群体未必完全相同，但它们共同指向一个事实：王室需要在西北方向长期投入军队、车马、粮食和封国协作。边疆不是地图边缘的一块空白，而是王室财政和军事能力每天都要面对的账本。

\begin{event}{宣王时期}{料民、北方战争与“中兴”叙事}{王年与细节有争议}{关中及西北}\eventanchor{WZ-XUAN-REVIVAL}{西周·宣王中兴}
\term{这一阶段发生了什么}：传世文献和《诗经》中的相关篇章保存了宣王时期对外用兵、整顿人口与王室试图恢复权威的记忆。它常被称为“宣王中兴”。

\term{事情为什么要谨慎看}：所谓中兴并不必然意味着王室重新回到成康时期的强度。恰恰相反，频繁料民和边境战争可能说明中心已经感到兵源、税源与控制力的不足。

\term{后来改变了什么}：宣王的努力或许延缓了危机，却没有消除封国坐大、边防耗费和继承风险。幽王时期的崩溃因此更像长时间压力累积后的断裂，而非一夜之间的道德惩罚。
\end{event}

\begin{textwindow}[《诗经》里的边防焦虑]
《诗经·小雅·采薇》写“靡室靡家，玁狁之故；不遑启居，玁狁之故”。它不是一份完整战报，却让我们听见长期戍役对家庭生活的压迫。诗歌提供的是情绪和社会感受，需与金文、史书和历史地理材料互相校验。
\end{textwindow}
